\documentclass[11pt]{article}

% --- Packages ---
\usepackage[margin=1in]{geometry}
\usepackage{amsmath, amssymb}
\usepackage{graphicx}
\usepackage{booktabs}
\usepackage{array}
\usepackage{hyperref}
\usepackage{natbib}
\usepackage{authblk}
\usepackage{caption}

\setlength{\parskip}{4pt}

% --- Title ---
\title{Factorized Visual Representations in the Primate Visual System and Deep Neural Networks}
\date{}

\begin{document}
\maketitle

\begin{abstract}
Classification of object identity has been proposed as a principal objective
of the primate ventral visual stream and is widely used as the training
target for deep neural network (DNN) models of that stream. But high-level
visual cortex simultaneously represents many other attributes --- pose,
position, lighting, viewpoint --- and classification alone does not say how
that non-identity information should be encoded. We argue here that a
quantifiable geometric property, factorization, is a normative principle
that constrains both biological and model representations. Factorization
asks whether variance driven by non-class scene parameters is placed in
orthogonal subspaces of population activity, rather than being discarded
(invariance) or mixed into the identity subspace (entangled coding). We
introduce complementary PCA-based and covariance-based measures of
factorization that apply to arbitrary scene variables. In macaque V4 and IT
multi-unit populations (\(n=128\) sites each), factorization of identity
from non-identity information increased from V4 to IT, and a representational
lesion that preserved mean firing rates, covariance, and invariance but
selectively reduced factorization significantly reduced object-class
decoding accuracy (chance \(=1/64\)). Across a diverse DNN library evaluated
on 12 primate datasets, greater factorization of four scene parameters
(background, lighting, object pose, camera viewpoint) predicted higher
neural, fMRI, and behavioral fit, while invariance predicted fit only for
appearance parameters. Combining factorization with classification in a
cross-validated regression (\(n=4\) datasets per target; 80/20 splits; 100
resamples) significantly improved predictions of the most brain-like models
over classification or dimensionality alone. Factorization thus emerges as a
widely used representational strategy in both brains and brain-like models.
\end{abstract}

\section{Introduction}

Recognizing objects across identity-preserving transformations is one of the
defining computational problems of primate vision \citep{dicarlo2007untangling,
dicarlo2012how}. Deep convolutional neural networks trained on large-scale
image classification have become the most predictive quantitative models of
spiking responses in macaque V4 and inferior temporal (IT) cortex
\citep{yamins2014performance, khalighrazavi2014deep, yamins2016using}. Across
many studies, models with higher classification accuracy tend to show higher
neural predictivity at both the single-unit and population level
\citep{majaj2015simple, schrimpf2020integrative}. Yet this relationship
saturates, and at the high end it sometimes reverses: beyond a threshold
level of classification performance, further gains do not buy better neural
predictions \citep{schrimpf2020integrative}. This plateau suggests that
classification alone is not a complete normative constraint on what it means
for a representation to be ``brain-like,'' and motivates searching for
additional principles that further structure the representation.

One natural candidate principle is invariance. A canonical way to support
linear readout of identity is to render the representation insensitive to
non-class image variation such as pose, background, or illumination
\citep{dicarlo2007untangling}. But invariance is intrinsically a zero-sum
strategy: information that is collapsed in the population code is
permanently unavailable for downstream use. Empirically, high-level ventral
stream neurons do not simply collapse non-identity information. They
represent object pose, position, and scale explicitly, and this explicit
encoding tends to grow from earlier to later stages of the hierarchy
\citep{hong2016explicit, rust2012balanced, freiwald2010functional}. Many
behaviors --- navigation, grasping, visually guided reaching --- require
simultaneous access to identity and to non-identity variables, which is in
tension with a purely invariant account.

An alternative is factorization: encoding different sources of image variance
in orthogonal, non-interfering subspaces of population activity. A factorized
population maintains the full dimensionality of the visual input while
reserving a specific subspace for linearly separable identity information.
This concept is closely related to, but more general than, the machine
learning notion of manifold disentanglement \citep{higgins2017beta}: it does
not require that independent generative factors project onto individual
neurons, only onto orthogonal subspaces. Factorization also connects to
recent work on abstraction and high-dimensional representational geometry in
prefrontal cortex, hippocampus, and trained neural networks
\citep{bernardi2020geometry, rigotti2013importance, fusi2016why,
stringer2019high, sorscher2022neural}. Whereas invariance asks how much
non-class variance remains, factorization asks where that variance lies
relative to the class-coding axes.

Here we ask whether factorization is empirically quantifiable, whether it
emerges in primate visual cortex, and whether it predicts which DNN models
best match neural and behavioral data. We introduce PCA-based and
covariance-based measures of factorization that apply to arbitrary scene
variables, including multi-dimensional generative factors such as object
pose and camera viewpoint. We apply these measures to existing macaque and
human datasets \citep{majaj2015simple, hong2016explicit, rajalingham2018large,
khalighrazavi2014deep}, to a diverse DNN library spanning supervised
\citep{krizhevsky2012imagenet, he2016deep}, contrastive self-supervised
\citep{chen2020simple, he2020momentum}, and auxiliary self-supervised
\citep{zhuang2021unsupervised} training objectives, and to natural-movie
datasets that approximate object-pose or camera-viewpoint variation
implicitly. We further introduce a representational lesion experiment: a
change-of-basis transformation that preserves mean firing rates, covariance,
and invariance while selectively reducing factorization, allowing us to ask
whether factorization causally supports decoding.

The paper makes four contributions. First, we formalize factorization as a
population-level representational property with two operationally
complementary definitions. Second, we show that factorization of object
identity from non-identity information grows from V4 to IT in macaque and
that reducing factorization through a representation-preserving rotation
significantly degrades object-class decoding (chance \(=1/64\); \(n=128\)
multi-unit sites in V4 and 128 in IT). Third, we demonstrate that across a
broad DNN library, factorization of each of four scene parameters
(background, lighting, object pose, camera viewpoint) is consistently
correlated with neural, fMRI, and behavioral predictivity across 12 primate
datasets, whereas invariance is predictive only for appearance parameters.
Fourth, we show that combining factorization with classification in a
cross-validated regression model --- averaged over \(n=4\) datasets for each
of V4, IT/HVC, and behavior, using 80/20 train/test splits and 100
resampling rounds --- substantially improves prediction of the most
brain-like models over classification or dimensionality alone. We conclude
that factorized encoding is a widely used representational strategy and is a
useful complement to classification performance as a criterion for
brain-like vision models.

\section{Results}

\subsection{A framework for quantifying factorization and invariance}

The problem of encoding object identity in a linearly separable way admits
three qualitatively different representational strategies
(Figure~\ref{fig:framework}). In a fully invariant code, non-identity
variance is collapsed into tight spheres within the identity subspace, so
non-identity information is discarded outright. In a factorized code,
non-identity variance is preserved along axes orthogonal to the identity
subspace, so that both identity and non-identity variables can be read out
linearly without mutual interference. In an entangled code, non-identity
variance overlaps the identity subspace, so that variation in a non-class
variable corrupts identity decoding. Factorization and invariance are not
mutually exclusive, but they make fundamentally different predictions about
how much non-class information a population preserves and where that
information lives.

We operationalize these ideas with two complementary metrics
(see~Methods). For a scene parameter of interest, the PCA-based factorization
score is defined as the fraction of parameter-induced variance that lies
outside the top principal components of the variance induced by all other
scene parameters, using a variance threshold of 90\% to define that
``other-parameter subspace.'' The covariance-based factorization score,
which has no free parameter, is defined as \(1\) minus the normalized inner
product of the covariance matrices induced by the parameter of interest and
by all other parameters. Invariance to a parameter is measured independently,
as the ratio between the variance induced by that parameter and the variance
induced by all other parameters. These two metrics give a factorization and
an invariance score for every (population, scene-parameter) pair.

To make the consequences of factorization concrete, we simulated a
two-variable coding problem in which \(N=10\) binary features jointly
determined a population response (Figure~\ref{fig:simulation}). Two of the
features were designated as decoding targets, and the angle between their
coding axes was varied from fully orthogonal (alignment \(a=0\), square
geometry) to fully aligned (\(a=1\), collinear). All other features were
encoded on mutually orthogonal axes, and isotropic Gaussian noise was added
at the neuron level. As alignment increased, a linear SVM trained on a
small number of samples showed progressively worse generalization on
held-out data, even though the two variables remained nominally linearly
separable. The simulation isolates the effect of factorization per se: in
the orthogonal regime, a reader that targets one variable automatically
averages out variance from the other, while in the entangled regime the two
variables bleed into each other and destroy the generalization benefit.

\subsection{Factorization increases from V4 to IT and supports object decoding}

Given the theoretical appeal of factorized representations, we next asked
whether primate visual cortex exhibits factorization and, if so, how much it
contributes to decoding. We used an existing macaque dataset in which a
grid of images independently varied object identity versus object pose and
background, recorded from multi-unit arrays in V4 and IT (\(n=128\) sites
per area) \citep{majaj2015simple, hong2016explicit}. To quantify
factorization, we compared the measured PCA-based factorization to a shuffle
control in which object identity labels were randomized across images; this
control removes factorization driven purely by changes in covariance
structure and dimensionality, so the residual is the signal we are after.

Both V4 and IT populations exhibited factorization scores significantly
above the shuffle control, and factorization was higher in IT than in V4
(Figure~\ref{fig:hierarchy}). Invariance to non-identity information also
increased from V4 to IT, consistent with earlier reports
\citep{rust2012balanced}. Importantly, when we measured invariance not over
the full population but within the subspace that captured the bulk (90\%)
of identity-driven variance, the V4-to-IT increase in invariance became even
more pronounced, reflecting the fact that increased factorization itself
lifts non-identity variance away from the identity axes and thus boosts
invariance specifically inside the identity subspace. These two properties
therefore reinforce each other in the empirical data, rather than competing.

To ask whether factorization actually matters for decoding --- not merely
whether it is present --- we designed a representational lesion experiment
targeted specifically at factorization (Figure~\ref{fig:hierarchy}). We
computed the principal components of the mean responses to each object class
(the inter-class PCs) and of the within-class residuals (the intra-class
PCs). We then applied a change-of-basis matrix \(W_{\text{inter}\to\text{intra}}\)
to the class centers that rotated each inter-class PC into the
corresponding intra-class PC. This transformation leaves intra-class
variance statistics exactly unchanged, and it also preserves the overall
variance induced by non-class factors and the invariance of the population
to those factors; what it specifically removes is the orthogonality of
inter-class variance with respect to intra-class variance, i.e.\ the
factorization of class from non-class variance. When we then trained a
linear decoder on the transformed population, object-class decoding accuracy
dropped significantly relative to the untransformed data in both V4 and IT
(chance \(=1/64\); \(n=128\) multi-unit sites per area). Invariance thus
does not, on its own, explain the decoding performance achievable from V4
and IT populations; factorization is required to recover the full accuracy.

\subsection{Factorization in DNN models predicts brain-likeness}

Having established that factorization is present in ventral stream data and
that it causally supports object-class decoding, we next turned to DNN
models. We had two motivations. First, DNN representations are the best
current quantitative models of IT cortex
\citep{yamins2014performance, yamins2016using, khalighrazavi2014deep},
so if factorization is a genuine organizing principle of biological vision
we should expect the most brain-like DNNs to be the most factorized. Second,
DNNs offer experimental flexibility that neural recordings do not: we can
independently vary any scene parameter and present thousands of stimuli,
which is what the factorization metric requires for a clean estimate. We
therefore asked whether DNN models exhibit factorization, whether it grows
during training, and whether across models it predicts how well a model
matches neural and behavioral data. DNN analyses offer a crucial advantage: because we can
present arbitrary stimuli, we can independently vary individual scene
parameters at sufficient scale to compute factorization and invariance for
each in isolation. We generated an augmented scene library based on 100
base scenes, in which each of four parameters --- foreground object pose,
background identity, scene lighting, and 2D camera viewpoint --- was varied
at \(n=10\) parameter levels while the others were held fixed
(Figure~\ref{fig:pipeline}). This yielded a 4000-image stimulus set per
model. Parameter ranges were taken from standard augmentation practices
used in self-supervised training pipelines
\citep{chen2020simple, he2020momentum}: object angle was varied from
\([-90,90]\) degrees along all three axes, in-plane object position from
\([-30\%,30\%]\) horizontally and \([-60\%,60\%]\) vertically, size from
\(\times 1/1.6\) to \(\times 1.6\), camera crop scale from 20\% to 100\% of
the image, and lighting jitters spanned \([-0.4,0.4]\) for brightness,
contrast, and saturation and \([-0.1,0.1]\) for hue.

We compared a diverse library of DNNs spanning supervised image
classification \citep{krizhevsky2012imagenet, he2016deep}, contrastive
self-supervised training \citep{chen2020simple, he2020momentum}, and
auxiliary self-supervised objectives \citep{zhuang2021unsupervised}. In
every case, training produced factorization scores above those of an
untrained, randomly initialized ResNet-50 control for all four scene
parameters. By contrast, training had mixed effects on invariance: it
typically increased invariance for background and lighting but reduced it
for object pose and camera viewpoint. Representational metrics computed in
a model's penultimate layer were positively correlated, across models, with
the model's ability to predict IT single-unit firing rates on the macaque E2
dataset, and separately with its ability to predict human per-image
behavioral signatures on the I2 matrix
\citep{rajalingham2018large}. A notable qualitative pattern emerged: high
invariance to object pose was associated with a drop in neural predictivity,
in contrast to the monotonic relationship observed for factorization.

\subsection{Generality across 12 datasets and measurement modalities}

A concern with any single comparison of this kind is that the pattern might
depend idiosyncratically on the stimulus set, brain area, or measurement
modality chosen. We therefore extended the analysis across 12 independent
datasets: 6 macaque sets (2 V4, 2 IT, 2 behavior) and 6 human sets (2 V4
fMRI, 2 HVC fMRI, 2 behavior)
\citep{majaj2015simple, hong2016explicit, rajalingham2018large,
khalighrazavi2014deep}. For each model, we averaged the factorization and
invariance scores over the final five representational layers, and
independently averaged neural or behavioral predictivity over the same
layers, to avoid idiosyncratic effects of any single layer.

Across all 12 datasets, factorization of each of the four scene parameters
was positively correlated with model brain-likeness. Invariance, by contrast,
split along an appearance-versus-spatial axis: invariance to background
identity and lighting was generally predictive of brain-likeness, while
invariance to object pose and camera viewpoint was not (zero or negative
correlation across datasets). This dissociation is striking because all
four parameters contribute similarly to how ``easy'' object classification
is: a model could, in principle, improve its classification accuracy by
becoming invariant to any of them. The fact that the brain's
representational fit does not reward invariance to spatial transforms
suggests that the ventral stream actively preserves spatial information
rather than discarding it, consistent with the direct neural observation of
explicit pose and position coding in IT \citep{hong2016explicit}.

We further probed the robustness of the finding along three methodological
axes. First, restricting the analysis to models sharing the ResNet-50
architecture preserved the factorization-vs-predictivity relationship,
ruling out architectural confounds. Second, replacing linear-encoding fits
with representational dissimilarity matrix (RDM) correlations
\citep{kriegeskorte2008representational, seeliger2017convolutional} left the
results qualitatively unchanged. Third, sweeping the PCA variance threshold
that defines the ``other-parameter subspace'' from 50\% to 99\% left the
main findings intact. We also confirmed that the covariance-based
factorization metric, which has no free parameters, gave similar results.
Finally, we recomputed factorization from natural movies rather than from
synthetic augmentations. Using frame-by-frame variance (contiguous frames
separated by 200\,ms) from the MIT Moments in Time movie set (stationary
observer, approximating object-motion variance) and the UT-Austin Egocentric
movie set (non-stationary observer, approximating camera-viewpoint
variance), factorization again predicted brain-likeness across macaque and
human datasets, while invariance did not. Thus the main finding --- that
factorization of scene parameters predicts brain-likeness --- holds across
brain signal (spiking vs.\ BOLD), species (monkey vs.\ human), cortical
areas (V4 vs.\ IT / HVC), and stimulus families (synthetic grayscale images
vs.\ natural color movies).

\subsection{Combining factorization with classification improves brain predictions}

A reasonable alternative explanation for the above results is that
factorization is simply a proxy for classification performance, which is
known to correlate with neural predictivity
\citep{yamins2014performance, majaj2015simple}. To test this, we examined
how classification-alone predictions of brain-likeness behaved at the
high-performance end of the DNN library. For IT and HVC datasets, the
relationship between model classification accuracy and neural or voxel
predictivity showed the saturating and sometimes reversing trend noted in
prior work \citep{schrimpf2020integrative}, consistent with the intuition
that beyond a certain accuracy the additional gains do not translate into
more brain-like features.

We then fit linear regressions in which both classification and
factorization (for a chosen scene parameter) were used as predictors of
neural or behavioral fit. The regression was trained on 80\% of models and
tested on the held-out 20\%, and this procedure was repeated over 100
random train/test splits to obtain a robust estimate. Results were averaged
across \(n=4\) datasets for each of V4, IT/HVC, and behavior. Linearly
combining factorization with classification produced significant
improvements in predicting the most brain-like models relative to using
classification alone, and the saturating trend on IT/HVC data was no longer
present when object-pose factorization was added as a second predictor. The
same boost was not observed when we combined classification with overall
representational dimensionality: classification plus dimensionality did no
better than classification alone on neural and fMRI data, showing that the
information carried by factorization is not reducible to the size of the
representational space. These effects were replicated with both the PCA-based
and covariance-based factorization metrics.

\section{Methods}

\subsection{Neural and behavioral datasets}

Analyses used twelve previously published datasets spanning macaque and human
primate visual cortex and behavior
\citep{majaj2015simple, hong2016explicit, rajalingham2018large,
khalighrazavi2014deep}. Macaque electrophysiology datasets included
single-unit recordings in V4 and anterior ventral IT, and chronic multi-unit
recordings across V4 and IT (the latter biased toward CIT and PIT, where
multi-unit arrays are more easily placed). Human fMRI datasets included two
color-image and two grayscale-image sets of V4 and HVC voxel responses at
comparable spatial resolution. Behavioral datasets provided image-by-image
object-classification performance for macaques and humans. A schematic
organization of these datasets is given in Table~\ref{tab:datasets}.

\begin{table}[htbp]
\centering
\caption{Datasets used for measuring DNN brain-likeness.}
\label{tab:datasets}
\begin{tabular}{lllc}
\toprule
Dataset category & Species & Region / modality & Count \\
\midrule
V4 single/multi-unit & Macaque & V4 cortical spiking & 2 \\
IT single/multi-unit & Macaque & IT cortical spiking & 2 \\
Macaque behavior & Macaque & Object classification (I1/I2) & 2 \\
V4 fMRI & Human & V4 BOLD & 2 \\
HVC fMRI & Human & High-level visual cortex BOLD & 2 \\
Human behavior & Human & Object classification (I1/I2) & 2 \\
\bottomrule
\end{tabular}
\end{table}

\subsection{DNN model library}

We tested DNN models drawn from three training-objective families, as
summarized in Table~\ref{tab:dnns}. Supervised models were trained on
large-scale image classification \citep{krizhevsky2012imagenet, he2016deep}.
Contrastive self-supervised models were trained with instance-discrimination
objectives \citep{chen2020simple, he2020momentum}. Auxiliary self-supervised
models were trained with non-contrastive unsupervised objectives such as
reconstruction or colorization \citep{zhuang2021unsupervised}. A majority of
the models shared the ResNet-50 architecture \citep{he2016deep}; a randomly
initialized ResNet-50 served as the untrained control in
Figure~\ref{fig:pipeline}.

\begin{table}[htbp]
\centering
\caption{DNN model families sampled for the main analyses.}
\label{tab:dnns}
\begin{tabular}{lll}
\toprule
Family & Training objective & Representative members \\
\midrule
Supervised & Image classification & AlexNet, ResNet-50 \\
Contrastive SSL & Instance discrimination & MoCo, SimCLR \\
Auxiliary SSL & Reconstruction / colorization / prediction & Unsupervised ensemble \\
Untrained control & None (random initialization) & ResNet-50 (random init) \\
\bottomrule
\end{tabular}
\end{table}

\subsection{Scene parameter augmentation}

For each of 100 base scenes drawn from the existing neural experimental
dataset, we generated sets of images varying exactly one of four scene
parameters --- foreground object pose, background identity, scene lighting,
and 2D camera viewpoint --- at \(n=10\) parameter levels, while holding all
other parameters fixed. This produced 4000 stimuli per model. Parameter
ranges are summarized in Table~\ref{tab:augparams} and follow standard
practice for self-supervised augmentation pipelines
\citep{chen2020simple, he2020momentum}.

\begin{table}[htbp]
\centering
\caption{Scene parameter augmentation ranges.}
\label{tab:augparams}
\begin{tabular}{ll}
\toprule
Parameter & Range \\
\midrule
Object angle (each of 3 axes) & [-90, 90] degrees \\
Object in-plane horizontal position & [-30\%, 30\%] \\
Object in-plane vertical position & [-60\%, 60\%] \\
Object size & [$\times 1/1.6$, $\times 1.6$] \\
Camera crop scale & 20\%--100\% of image \\
Brightness / contrast / saturation jitter & [-0.4, 0.4] \\
Hue jitter & [-0.1, 0.1] \\
\bottomrule
\end{tabular}
\end{table}

\subsection{Factorization and invariance metrics}

For a chosen scene parameter \(p\), let \(\mathrm{var}_{p}\) denote the
total variance (summed across all neural dimensions) of population responses
as \(p\) is varied across its 10 levels, averaged over the 100 base scenes,
and let \(\mathrm{var}_{\mathrm{other}}\) denote the corresponding quantity
when all other scene parameters are varied. Invariance to \(p\) is defined
as a monotonically decreasing function of
\(\mathrm{var}_{p}/\mathrm{var}_{\mathrm{other}}\): the smaller the variance
induced by \(p\) relative to the variance induced by all other parameters,
the higher the invariance to \(p\). Higher values correspond to neural
responses that are less perturbed by variation in \(p\).

\paragraph{PCA-based factorization.}
The ``other-parameter subspace'' was defined by averaging neural responses
for each base scene over all augmentations of the parameter of interest and
running PCA on the resulting set. The subspace was taken as the span of the
top principal components containing 90\% of the variance of these responses.
Let \(\mathrm{var}_{p\mid\mathrm{other}}\) denote the variance of
\(p\)-induced activity inside that subspace. Factorization is then defined
as the fraction of \(p\)-induced variance that lies outside it:
\begin{equation}
\mathrm{Fact}_p^{\mathrm{PCA}}
= 1 - \frac{\mathrm{var}_{p\mid\mathrm{other}}}{\mathrm{var}_{p}}.
\end{equation}

\paragraph{Covariance-based factorization.}
Let \(\mathrm{cov}_{p}\) and \(\mathrm{cov}_{\mathrm{other}}\) be the
covariance matrices of \(p\)-induced and other-parameter-induced activity,
respectively. The covariance-based factorization score is defined as
\begin{equation}
\mathrm{Fact}_p^{\mathrm{cov}}
= 1 - \frac{\langle \mathrm{cov}_{p},\ \mathrm{cov}_{\mathrm{other}}\rangle}
{\lVert \mathrm{cov}_{p}\rVert\ \lVert \mathrm{cov}_{\mathrm{other}}\rVert},
\end{equation}
where \(\langle\cdot,\cdot\rangle\) denotes the elementwise inner product of
flattened matrices and \(\lVert\cdot\rVert\) denotes the Frobenius norm.
This measure has no free parameters. Across all subsequent analyses, the
PCA-based and covariance-based metrics produced similar findings, supporting
the robustness of the factorization concept.

\subsection{Representational lesion transformation}

To test the necessity of factorization for decoding, we applied a change-of-
basis transformation that selectively removed orthogonality between
inter-class and intra-class variance axes. Letting \(x_c\) be the mean
neural activity vector for class \(c\), we computed the principal components
of the class centers (the inter-class PCs) and the principal components of
the residuals \(x - x_c\) (the intra-class PCs). We then constructed a
rotation matrix \(W_{\text{inter}\to\text{intra}}\) that aligned each
inter-class PC with the corresponding intra-class PC, and applied this
rotation to the class centers only, leaving within-class activity
statistics untouched. Explicitly, for an activation \(x\) in class \(c\),
the transformed activation was
\(x' = W_{\text{inter}\to\text{intra}}\,x_c + (x - x_c)\). By construction,
this transformation preserved (i) mean firing rates per class, (ii) the
covariance structure of intra-class variance, and (iii) the invariance of
the population to non-class factors, while specifically reducing
factorization of class from non-class variance.

\subsection{Natural-movie factorization}

For natural movies, variance was defined implicitly by the temporal
structure of the stimulus rather than by parametric control. We used the MIT
Moments in Time movie set, in which the observer is largely stationary and
the dominant source of temporal change is object and scene motion, and the
UT-Austin Egocentric movie set, in which an observer translates extensively
in an urban environment and the dominant change is camera viewpoint.
Contiguous frames separated by 200\,ms were treated as reflecting local
variance in one of the two motion parameters; variance across non-contiguous
frames provided the ``all parameters'' baseline.

\subsection{Neural-encoding fits and cross-validated combination}

Linear mappings from model features to neural or voxel responses were fit
using ridge regression with regularization selected by cross-validation,
applied to the top 300 principal components of model features. As an
alternative, RDM-based comparisons were computed using dot-product
similarities of image pairs and Spearman correlation between model and
neural RDMs \citep{kriegeskorte2008representational}. For the combined
metric analyses, a linear regression was fit using both classification
accuracy and a chosen factorization score as predictors of neural or
behavioral fit; the regression was trained on 80\% of models and tested on
the held-out 20\%, with 100 random train/test splits. Results were averaged
across \(n=4\) datasets for each of V4, IT/HVC, and behavior. Correlation
between predicted and actual neural or behavioral fit on held-out models
was reported as Spearman correlation.

\section{Discussion}

The central finding of this work is that factorization --- the degree to
which different scene variables are encoded in orthogonal subspaces of a
population code --- is a quantitative signature shared by high-level
primate visual cortex and by the deep neural network models that most
closely predict it. Factorization is not reducible to classification
performance or to representational dimensionality, and it carries
information about brain-likeness that these more established properties
miss. We present this as evidence that factorization is a normative
principle of primate visual representations, complementary to classification
performance and invariance. In the macaque ventral stream, factorization of object
identity from non-identity information grows from V4 to IT and is
causally required to support the full object-class decoding accuracy
observed in these populations. A representational lesion that preserved
mean firing rates, covariance, and invariance while selectively reducing
factorization significantly degraded decoding, isolating factorization as
the operative geometric property rather than as a correlate of more
familiar statistics. Across a diverse library of DNN models evaluated on
12 primate datasets, factorization of four scene parameters consistently
predicted how brain-like a model is, while invariance did so only for
appearance parameters. These results do not contradict the established role
of classification \citep{yamins2014performance, majaj2015simple}; rather,
they complement it. A cross-validated regression combining factorization
with classification produced significant improvements in predicting the
most brain-like models, and these improvements were not obtained by
combining classification with dimensionality alone.

Our definition of factorization occupies a specific place in a growing
landscape of ideas about population geometry. Disentangled representation
learning in machine learning seeks to recover independent generative
factors as individual latent axes \citep{higgins2017beta}; factorization as
defined here relaxes that single-axis assumption to the case of orthogonal
\emph{subspaces}, which is a more natural description of how
high-dimensional generative factors such as pose and viewpoint are encoded.
Abstraction and shattering dimensionality in prefrontal cortex and
hippocampus \citep{bernardi2020geometry, rigotti2013importance, fusi2016why}
are closely related: a fully factorized code should be highly abstract, in
the sense that linear decoders for one variable generalize across values of
other variables. Our emphasis is more geometric --- we directly characterize
the arrangement of subspaces --- while abstraction-based analyses emphasize
the downstream decoding consequences. These are two views of a single
underlying structure, and our lesion result makes explicit the link between
the geometric and the decoding-based accounts.

Factorization is also distinct from, and complementary to, representational
straightening \citep{henaff2019perceptual, henaff2021primary}. Straightening
concerns the local temporal geometry of trajectories \emph{within} a
variable's subspace --- how curved they are as a function of time in the
population space --- while factorization concerns the global arrangement of
those subspaces with respect to one another. Straightening has been shown
to matter for early visual areas, where predictive extrapolation of input is
a candidate goal, while factorization emerges more prominently in higher
visual cortex, where simultaneously readable multi-variable codes are more
useful. The two concepts are mutually compatible and together form a richer
description of visual population geometry than either alone. A related line
of work has linked the shape of object manifolds to few-shot classification
performance \citep{sorscher2022neural}; high-dimensional but orthogonal
organization of non-class variance, as we define it, is exactly the kind of
structure that facilitates such generalization, and extending our
factorization framework to incorporate manifold-shape statistics would be a
natural next step.

Several limitations of our study are worth stating explicitly. First, our
analysis is correlational at the DNN level: we observed that models with
higher factorization tend to be more brain-like, but we did not train
networks with factorization as an explicit loss term. A causal test would
require direct optimization of factorization in a model's objective
function, and comparison of the resulting model to both factorization-blind
controls and to neural data. Second, the four scene parameters we examined
were chosen heuristically on the basis of what a standard augmentation
pipeline controls. Unsupervised discovery of the latent factors that
actually dominate the statistics of natural images --- for instance, using
methods related to independent component analysis --- would be a more
principled choice, and could reveal factorization structure that our
parameter-specific analyses miss. Third, our metrics describe the
macroscopic arrangement of subspaces but leave open the microscopic
geometry of coding \emph{within} each subspace. A population could achieve a
given factorization score with very different single-neuron tuning
profiles: each neuron might encode a single latent variable, distinct
sub-regions might encode qualitatively different variables, or each neuron
might carry a distributed mixed code
\citep{rigotti2013importance, fusi2016why, freiwald2010functional,
hong2016explicit}. Disambiguating these possibilities requires systematic
single-neuron and subregion-level analyses that go beyond the population
metrics used here.

Looking forward, factorization suggests concrete paths for building more
brain-like models. Training objectives that explicitly reward orthogonality
of parameter-induced subspaces, or that preserve parameter information in
non-interfering directions, could push models toward representations that
are simultaneously good at object classification and good at predicting
neural responses --- a combination that, as our cross-validated analyses
show, is captured better by classification plus factorization than by
classification alone. Finally, because the factorization metrics are
straightforward to compute on any neural or model representation, they
provide a tool for probing factorization in brain areas beyond visual
cortex. We expect that factorization will prove to be a useful lens not
just for vision but for any system in which multiple behaviorally relevant
variables must be simultaneously read out from a shared population code.

\begin{figure}[htbp]
\centering
\includegraphics[width=0.92\textwidth]{figures/fig_framework.pdf}
\caption{Three representational strategies for encoding object identity in
the presence of non-identity variance. Factorization places non-identity
variance along axes orthogonal to the identity subspace (orange line).
Invariance collapses non-identity variance into compact regions. Entangled
codes allow non-identity variance to span directions overlapping the
identity subspace, corrupting decoding.}
\label{fig:framework}
\end{figure}

\begin{figure}[htbp]
\centering
\includegraphics[width=0.82\textwidth]{figures/fig_simulation.pdf}
\caption{Simulation schematic for two binary variables encoded out of
\(N=10\) total dimensions. Orthogonal codes (alignment \(a=0\)) support
accurate linear decoding of both variables from few samples under noise;
non-orthogonal codes (\(a>0\)) preserve separability but degrade
generalization.}
\label{fig:simulation}
\end{figure}

\begin{figure}[htbp]
\centering
\includegraphics[width=0.92\textwidth]{figures/fig_hierarchy.pdf}
\caption{Macaque ventral visual hierarchy and the empirical trend reported
here. Factorization and invariance both increase from V4 to IT
(\(n=128\) multi-unit sites per area). A change-of-basis lesion that
preserves firing rates, covariance, and invariance but reduces factorization
significantly decreases object-class decoding (chance \(=1/64\)).}
\label{fig:hierarchy}
\end{figure}

\begin{figure}[htbp]
\centering
\includegraphics[width=0.98\textwidth]{figures/fig_pipeline.pdf}
\caption{DNN evaluation pipeline. For each of 100 base scenes, 10 parameter
levels of each of four scene parameters are generated, yielding 4000 images
per model. Factorization and invariance are computed from model activations
(PCA-based or covariance-based), and model predictivity is evaluated
against 12 primate datasets spanning V4, IT/HVC, and behavior.}
\label{fig:pipeline}
\end{figure}

\bibliographystyle{plainnat}
\bibliography{references}

\end{document}
